% Generated draft: requires mathematical and source review.
\documentclass[11pt]{article}
\usepackage[margin=25mm]{geometry}
\usepackage{fontspec}
\setmainfont{Latin Modern Roman}
\usepackage{amsmath,amssymb,mathtools,unicode-math}
\setmathfont{Latin Modern Math}
\usepackage{enumitem,needspace,xurl,hyperref}
\setlength{\parindent}{0pt}
\setlength{\parskip}{5pt}
\setlength{\emergencystretch}{3em}
\newcommand{\N}{\mathbb{N}}
\newcommand{\Z}{\mathbb{Z}}
\newcommand{\Q}{\mathbb{Q}}
\newcommand{\R}{\mathbb{R}}
\newcommand{\C}{\mathbb{C}}
\newcommand{\F}{\mathbb{F}}
\newcommand{\A}{\mathbb{A}}
\newcommand{\PP}{\mathbb{P}}
\newcommand{\E}{\mathbb{E}}
\newcommand{\eps}{\varepsilon}
\newcommand{\dd}{\,\mathrm{d}}
\begin{document}
\section*{UnsolvedMath ID: 30005964}
\textbf{Source problem code:} catalogue entry
\subsection{Definitions and mathematical statement}
\textbf{Source formulation (Conjecture 6; incomplete setup).} Let \(S\), \(v\), the pairing \(\langle\ ,\ \rangle\), and the reduced category \(T_S(v)_w^{\mathrm{red}}\) have the meanings intended in the source report. Those meanings and their required hypotheses are not specified in the supplied record. In particular, the phrase ``\(w\) coprime with \(v\)'' is retained as a source condition; it is not defined here because the nature of \(v\) is unspecified.

The proposed assertion is that, for every \(g\geq 0\) in the source's intended parameter domain and every \(w\in\Z\) coprime with \(v\) in the source's sense, setting
\[
 g'=\frac{\langle v,v\rangle+2}{2},
\]
one has
\[
 T_S(v)_w^{\mathrm{red}}
 \quad\text{deformation equivalent to}\quad
 D^b\!\left(\operatorname{Hilb}(S,g')\right).
\]
Here the conventional interpretation of \(\operatorname{Hilb}(S,n)\) is the Hilbert scheme of zero-dimensional subschemes of \(S\) of length \(n\), and that of \(D^b(X)\) is the bounded derived category of coherent sheaves on \(X\). These interpretations require confirmation against the source's setup. In particular, the record supplies no hypotheses ensuring that \(g'\) is a nonnegative integer. It also does not specify the definition of deformation equivalence for the categories in this assertion, the construction denoted by \(T\), or the operation denoted by \(\mathrm{red}\). The parameter \(g\) does not otherwise occur in the displayed assertion; its domain beyond \(g\geq0\), and its relation to \(S\) or \(v\), are missing. Consequently, this is a transcription of the intended conjecture with unresolved qualifications, rather than a fully specified standalone mathematical statement.

\textbf{Qualifications retained from the surrounding extract.} The source expresses an expectation that the categories written there as \(T_S^\sigma(v)_w\) do not produce new examples of irreducible holomorphic symplectic varieties, and a belief that they are related to known examples by interesting derived equivalences, more specifically by compactifications of the Fourier--Mukai derived equivalence between a possibly relative abelian variety and its dual. It promises a more precise conjecture in future work and describes Conjecture 6 as a simplified form and a categorical version of a conjecture numbered (4), whose statement is absent.

The supplied original extract also contains adjacent clauses labelled (ii) and (iii). Clause (ii) says that the Serre functor \(S_T\) of \(T_S^\sigma(v)_w^{\mathrm{red}}\) is trivial \(\acute{e}\)tale locally on \(M_S^\sigma(v)\). Clause (iii) asserts a wall-crossing equivalence between reduced categories associated with stability parameters, but its prime and superscript placement is corrupted in the supplied text. The definitions of \(\sigma\), \(M_S^\sigma(v)\), and local triviality of the Serre functor are absent, as are the hypotheses and precise indexing of the wall-crossing claim. These clauses are preserved as surrounding source claims; the extract does not establish that they are additional conclusions of Conjecture 6.
\subsection{Short English statement}
Conjecture 6 proposes that each admissible reduced category \(T_S(v)_w^{\mathrm{red}}\), with \(w\in\Z\) coprime with \(v\) in the source's sense, deforms to the bounded derived category of the Hilbert scheme of \((\langle v,v\rangle+2)/2\) points on \(S\). The supplied record omits the setup needed to define admissibility, the categories, and deformation equivalence, and does not explain the role of \(g\geq0\). Its open status is supplied dated curation, not an independently verified conclusion.
\subsection{Sources}
\begin{itemize}
\item Local UnsolvedMath registry: supplied complete record, permanent ID \(30005964\), exact source problem code \texttt{catalogue entry}, title ``Deformation Equivalence of Categorical Hyperk\"ahler Constructions,'' in the collection ``Oberwolfach Reports Open Problems.'' No individual proposer is supplied.
\item Supplied report citation: \emph{Algebraic Geometry: Wall Crossing and Moduli Spaces, Varieties and Derived Categories} (2025), Oberwolfach Reports, DOI \texttt{10.4171/owr/2024/32}, pp.\ 1787--1854. Supplied source URL: \url{https://doi.org/10.4171/owr/2024/32}. The extract labels the principal assertion Conjecture 6. The registry separately records proposed year 2024.
\item Supplied literature reference, labelled ``original OWR report'': \url{https://doi.org/10.4171/owr/2024/32}.
\item Literature assessment dated exactly \texttt{2026-08-22}: the supplied curation labels the problem open, says the proposed deformation-equivalence statement is explicitly labelled Conjecture 6, and reports that no primary proof or disproof was located. This is supplied dated curation, not independently verified literature research. The background expressly requests independent verification of time-sensitive status and cited claims.
\item The registry labels the statement \texttt{corrected\_verified} and says it was reconstructed and checked against the official Oberwolfach report. These are supplied verification claims; no independent source inspection or mathematical verification is asserted in this draft.
\end{itemize}
\end{document}
